% !TEX program = xelatex
% Lernplatt - by Can Siebert
% Kurs 22 (Bo 143) - Tag 14 - Aufgabenheft
% Chatbots, Conversational Commerce & digitale Vertriebstools

\documentclass[11pt,a4paper,oneside]{article}

\usepackage{polyglossia}
\setdefaultlanguage{german}

\usepackage[a4paper,margin=2.5cm]{geometry}

\usepackage{fontspec}
\defaultfontfeatures{Ligatures=TeX}

\IfFontExistsTF{Charter}{%
  \setmainfont{Charter}%
}{%
  \IfFontExistsTF{Bitstream Charter}{%
    \setmainfont{Bitstream Charter}%
  }{%
    \setmainfont{TeX Gyre Schola}%
  }%
}

\IfFontExistsTF{Source Sans 3}{%
  \setsansfont{Source Sans 3}%
}{%
  \IfFontExistsTF{Source Sans Pro}{%
    \setsansfont{Source Sans Pro}%
  }{%
    \setsansfont{TeX Gyre Heros}%
  }%
}

\newcommand{\caveatfontfallback}{\itshape}
\IfFontExistsTF{Caveat}{%
  \newfontfamily\caveatfont{Caveat}[Scale=1.15]%
}{%
  \newcommand\caveatfont{\caveatfontfallback}%
}

\setmonofont{Latin Modern Mono}

\usepackage{microtype}
\linespread{1.15}

\usepackage{xcolor}
\definecolor{lpInk}{RGB}{20,28,38}
\definecolor{lpAccent}{RGB}{22,90,140}
\definecolor{lpAccentLight}{RGB}{210,228,240}
\definecolor{lpAmber}{RGB}{222,138,30}
\definecolor{lpAmberLight}{RGB}{255,243,217}
\definecolor{lpAmberBorder}{RGB}{196,118,18}
\definecolor{lpGray}{RGB}{96,104,114}
\definecolor{lpGrayLight}{RGB}{238,240,243}
\definecolor{lpGreen}{RGB}{34,120,72}
\definecolor{lpGreenLight}{RGB}{222,238,228}
\definecolor{lpGreenBorder}{RGB}{24,96,58}

\definecolor{lpL1}{RGB}{34,120,72}
\definecolor{lpL2}{RGB}{22,90,140}
\definecolor{lpL3}{RGB}{170,42,52}

\usepackage[most]{tcolorbox}
\tcbuselibrary{breakable,skins}

\newtcolorbox{infobox}[1][Hinweis]{%
  enhanced, breakable,
  colback=lpAccentLight,
  colframe=lpAccent,
  boxrule=0.6pt,
  arc=2pt,
  left=10pt,right=10pt,top=8pt,bottom=8pt,
  fonttitle=\sffamily\bfseries\color{white},
  coltitle=white,
  title={#1},
  attach boxed title to top left={xshift=8pt,yshift=-8pt},
  boxed title style={size=small, colback=lpAccent, colframe=lpAccent, boxrule=0.4pt, arc=1pt},
  top=14pt
}

\newtcolorbox{antwortbox}[1][Ihre Antwort]{%
  enhanced, breakable,
  colback=white,
  colframe=lpGray,
  boxrule=0.4pt,
  arc=2pt,
  left=10pt,right=10pt,top=8pt,bottom=8pt,
  fonttitle=\sffamily\bfseries\color{white},
  coltitle=white,
  title={#1},
  attach boxed title to top left={xshift=8pt,yshift=-8pt},
  boxed title style={size=small, colback=lpGray, colframe=lpGray, boxrule=0.4pt, arc=1pt},
  top=14pt
}

\usepackage{enumitem}
\setlist{itemsep=2pt, topsep=4pt, parsep=0pt}
\setlist[itemize,1]{label={\textbullet},leftmargin=1.6em}
\setlist[itemize,2]{label={\textendash},leftmargin=1.4em}
\setlist[enumerate,1]{label=\arabic*., leftmargin=1.8em}

\usepackage{tabularx}
\usepackage{array}
\usepackage{booktabs}
\usepackage{longtable}

\usepackage{fancyhdr}
\usepackage{lastpage}
\pagestyle{fancy}
\fancyhf{}
\renewcommand{\headrulewidth}{0.3pt}
\renewcommand{\footrulewidth}{0pt}
\fancyhead[L]{\footnotesize\sffamily\color{lpGray}Aufgabenheft \textperiodcentered{} Kurs 22 \textperiodcentered{} Tag 14}
\fancyhead[R]{\footnotesize\sffamily\color{lpGray}Chatbots \& Conversational Commerce}
\fancyfoot[C]{\footnotesize\sffamily\color{lpGray}\thepage{} / \pageref{LastPage}}

\fancypagestyle{plain}{%
  \fancyhf{}%
  \renewcommand{\headrulewidth}{0pt}%
  \fancyfoot[C]{\footnotesize\sffamily\color{lpGray}\thepage{} / \pageref{LastPage}}%
}

\usepackage[explicit]{titlesec}
\titleformat{\section}[block]
  {\sffamily\Large\bfseries\color{lpAccent}}
  {\thesection.}{0.6em}{#1}
  [{\vspace{2pt}\color{lpAccent}\titlerule[0.6pt]}]
\titleformat{\subsection}[block]
  {\sffamily\large\bfseries\color{lpInk}}
  {\thesubsection}{0.6em}{#1}
\titlespacing*{\section}{0pt}{14pt}{8pt}
\titlespacing*{\subsection}{0pt}{10pt}{4pt}

\usepackage[hidelinks,unicode]{hyperref}

\newcommand{\akzent}[1]{{\caveatfont\color{lpAmber}#1}}
\newcommand{\fachbegriff}[1]{\textbf{#1}}

\newcommand{\niveaubadge}[2]{%
  \tcbox[on line, colback=#1, colframe=#1, coltext=white,
         boxrule=0pt, arc=2pt, left=5pt,right=5pt,top=1pt,bottom=1pt,
         nobeforeafter]{\sffamily\bfseries\footnotesize #2}%
}
\newcommand{\nivLi}{\niveaubadge{lpL1}{L1\,\textbullet\,Wissen}}
\newcommand{\nivLii}{\niveaubadge{lpL2}{L2\,\textbullet\,Anwendung}}
\newcommand{\nivLiii}{\niveaubadge{lpL3}{L3\,\textbullet\,Management}}

\newcommand{\zeit}[1]{%
  \tcbox[on line, colback=lpGrayLight, colframe=lpGrayLight, coltext=lpInk,
         boxrule=0pt, arc=2pt, left=5pt,right=5pt,top=1pt,bottom=1pt,
         nobeforeafter]{\sffamily\footnotesize\textbf{$\sim$\,#1}}%
}

\newcommand{\aufgabenkopf}[4]{%
  \par\vspace{6pt}
  \noindent{\sffamily\Large\bfseries\color{lpAccent}Aufgabe #1 \textendash{} #2}\par
  \vspace{2pt}
  \noindent{\color{lpAccent}\rule{\linewidth}{0.6pt}}\par
  \vspace{4pt}
  \noindent #3 \hfill \zeit{#4}\par
  \vspace{6pt}
}

\newcommand{\ytquery}[1]{%
  \par\vspace{4pt}
  \noindent{\footnotesize\sffamily\color{lpGray}\textbf{YouTube-Suchquery:} \texttt{#1}}\par
}

\newcommand{\antwortlinien}[1]{%
  \par\vspace{6pt}
  \foreach \x in {1,...,#1}{%
    \noindent\makebox[\linewidth]{\dotfill}\\[6pt]
  }
}
\usepackage{pgffor}

% =========================================================================
\begin{document}

\begin{titlepage}
  \pagestyle{empty}
  \vspace*{1.5cm}
  \noindent{\sffamily\footnotesize\color{lpGray}LERNPLATT \textperiodcentered{} BY CAN SIEBERT}\\[2pt]
  \noindent{\color{lpAccent}\rule{\linewidth}{1.2pt}}\\[4pt]
  \noindent{\sffamily\footnotesize\color{lpGray}AUFGABENHEFT \textperiodcentered{} MODUL 7 \textperiodcentered{} TAG 14 \textperiodcentered{} KURS 22 (Bo 143) \textperiodcentered{} 17.\,Juni 2026}

  \vspace{3cm}
  {\sffamily\fontsize{30}{34}\selectfont\bfseries\color{lpInk}Aufgabenheft\\Tag 14\par}

  \vspace{0.7cm}
  {\sffamily\large\color{lpGray}Chatbots, Conversational Commerce\\\& digitale Vertriebstools \textendash{}\\Kundeninteraktion verantwortlich\\automatisieren.\par}

  \vspace{1.5cm}
  {\caveatfont\LARGE\color{lpAmber}11 Aufgaben \textperiodcentered{} L1 \textperiodcentered{} L2 \textperiodcentered{} L3}\par

  \vfill

  \noindent\begin{minipage}{0.55\linewidth}
    {\sffamily\footnotesize\color{lpGray}Niveau-Mix:}\\
    {\sffamily\small\color{lpInk}3\,\textsf{L1} \textperiodcentered{} 5\,\textsf{L2} \textperiodcentered{} 3\,\textsf{L3}}\\[10pt]
    {\sffamily\footnotesize\color{lpGray}Bearbeitung:}\\
    {\sffamily\small\color{lpInk}Selbstbearbeitung im Lernpfad}
  \end{minipage}\hfill
  \begin{minipage}{0.4\linewidth}
    \raggedleft
    {\sffamily\footnotesize\color{lpGray}Dozent:}\\
    {\sffamily\small\color{lpInk}Can Siebert}\\[10pt]
    {\sffamily\footnotesize\color{lpGray}Begleitend:}\\
    {\sffamily\small\color{lpInk}Lehrskript \& Lösungsheft}
  \end{minipage}

  \vspace{0.5cm}
  \noindent{\color{lpAccent}\rule{\linewidth}{0.6pt}}
\end{titlepage}

\section*{So arbeiten Sie mit diesem Heft}
\thispagestyle{plain}

\noindent Dieses Aufgabenheft enthält elf Aufgaben zu Tag 14 \textendash{} \emph{Chatbots, Conversational Commerce und digitale Vertriebstools}. Im Mittelpunkt steht nicht die Technologie, sondern die Frage, wie automatisierte Kundeninteraktion Wert schafft, gemessen wird und verantwortlich begrenzt bleibt.

\begin{itemize}
  \item \nivLi{} \textendash{} \fachbegriff{Wissen.} Kurze Reproduktion und Einordnung. 5\,--\,10 Minuten pro Aufgabe.
  \item \nivLii{} \textendash{} \fachbegriff{Anwendung.} Transfer auf konkrete Vertriebs- und Servicesituationen. 15\,--\,30 Minuten.
  \item \nivLiii{} \textendash{} \fachbegriff{Management.} Entscheidungsarchitektur und Verantwortung. 30\,--\,45 Minuten.
\end{itemize}

\begin{infobox}[Hinweis]
Die Musterlösungen finden Sie im separaten \emph{Lösungsheft Tag 14}. Beim Lesen lohnt es sich, an jeder Stelle die Frage zu stellen: \emph{Wo verbessert Automatisierung Kundennutzen \textendash{} und wo ersetzt sie ihn nur?}
\end{infobox}

\clearpage

% =========================================================================
% L1
% =========================================================================
\section*{Niveau L1 \textendash{} Wissen \& Reproduktion}
\addcontentsline{toc}{section}{L1 -- Wissen}

% --- AUFGABE 1 -----------------------------------------------------------
% LERNPLATT_HILFSVIDEOS
\section*{Hilfsvideos zu diesem Tag}
\addcontentsline{toc}{section}{Hilfsvideos zu diesem Tag}
\thispagestyle{plain}

\noindent Die folgenden YouTube-Erkl\"arvideos vermitteln die Inhalte, die Sie zur Bearbeitung der Aufgaben ben\"otigen. Klicken Sie auf den Titel, um das Video direkt zu \"offnen; die vollst\"andige URL steht in Klammern.

\begin{description}[style=nextline,leftmargin=18pt,labelindent=0pt,itemsep=8pt]
  \item[1.~Chatbots einfach erklärt — Grundlagen automatisierter Kundeninteraktion] \href{https://youtu.be/3ya1k29C3bA}{\textcolor{lpAccent}{https://youtu.be/3ya1k29C3bA}}\\[2pt]
  {\footnotesize\color{lpGray}(URL: \nolinkurl{https://youtu.be/3ya1k29C3bA})}
  \item[2.~Live-Chat und KI-Widget für Kundenservice und Lead-Gewinnung] \href{https://youtu.be/_mCADLc16yg}{\textcolor{lpAccent}{https://youtu.be/_mCADLc16yg}}\\[2pt]
  {\footnotesize\color{lpGray}(URL: \nolinkurl{https://youtu.be/_mCADLc16yg})}
  \item[3.~Wann lohnt sich ein Chatbot für kleine Unternehmen?] \href{https://youtu.be/4ks3RY6lSRQ}{\textcolor{lpAccent}{https://youtu.be/4ks3RY6lSRQ}}\\[2pt]
  {\footnotesize\color{lpGray}(URL: \nolinkurl{https://youtu.be/4ks3RY6lSRQ})}
  \item[4.~KI-Chatbot für den Kundenservice — Schritt-für-Schritt] \href{https://youtu.be/0beblzCWAF8}{\textcolor{lpAccent}{https://youtu.be/0beblzCWAF8}}\\[2pt]
  {\footnotesize\color{lpGray}(URL: \nolinkurl{https://youtu.be/0beblzCWAF8})}
\end{description}
\vspace{0.6em}
\noindent{\footnotesize\color{lpGray}\textit{Tipp:} \"Offnen Sie das Video, lassen Sie es im Hintergrund laufen und arbeiten Sie parallel an der Aufgabe.}

\clearpage

\aufgabenkopf{1}{Chatbot-Typen und Conversational Commerce}{\nivLi}{8\,min}

\noindent Definieren Sie in jeweils einem Satz:

\begin{enumerate}
  \item \textbf{Regelbasierter Chatbot} (Entscheidungsbaum, klar definierte Pfade).
  \item \textbf{KI-gestützter Chatbot} (Sprachmodell, freie Eingaben, generative Antworten).
  \item \textbf{Hybrider Bot-Mensch-Prozess} (Bot mit Eskalation an einen Menschen).
  \item \textbf{Conversational Commerce} (vertriebsorientierte Kommunikation über Chat, Messenger, Portal).
\end{enumerate}

\noindent Ergänzen Sie einen abschließenden Satz: Welche \emph{drei} Faktoren entscheiden, welcher Typ in einer konkreten Vertriebssituation der richtige ist?

\ytquery{Chatbot Typen regelbasiert KI hybrid Conversational Commerce Definition}

\antwortlinien{10}

\clearpage

% --- AUFGABE 2 -----------------------------------------------------------
\aufgabenkopf{2}{Einsatzfelder im B2B-Vertrieb}{\nivLi}{8\,min}

\noindent Acht typische Einsatzfelder für Chatbots und Conversational-Commerce-Tools im B2B-Vertrieb:

\begin{enumerate}
  \item Erstinformation und FAQ
  \item Leadqualifizierung
  \item Produktberatung
  \item Terminbuchung
  \item Angebotsvorbereitung
  \item Statusabfragen (Bestellung, Lieferung)
  \item Nachbestellungen / Reorder
  \item After-Sales-Service \& Reaktivierung
\end{enumerate}

\noindent Beschreiben Sie für jedes Einsatzfeld in einem Halbsatz: Welches \emph{Kundenproblem} wird gelöst? Welches \emph{Vertriebsproblem} wird gelöst?

\ytquery{Chatbot B2B Vertrieb Einsatzfelder Leadqualifizierung FAQ Nachbestellung}

\antwortlinien{14}

\clearpage

% --- AUFGABE 3 -----------------------------------------------------------
\aufgabenkopf{3}{ROI-Logik und KPIs}{\nivLi}{8\,min}

\noindent Sechs typische KPIs für die Bewertung von Chatbots und digitalen Vertriebstools:

\begin{enumerate}
  \item \textbf{Bot-Nutzungsrate} (wie viele Besucher / Kunden nutzen den Bot?)
  \item \textbf{Lösungsquote} (Anteil der Anfragen, die der Bot ohne Eskalation klärt)
  \item \textbf{Eskalationsquote} (Anteil, der an Menschen übergeben wird)
  \item \textbf{Leadqualifizierungsrate} (Anteil qualifizierter Leads aus Bot-Dialogen)
  \item \textbf{Reaktionszeit} (durchschnittliche Antwortzeit)
  \item \textbf{Kosten pro Kontakt} (vor / nach Bot-Einführung)
\end{enumerate}

\noindent Ergänzen Sie eine kompakte \emph{ROI-Formel} in einem Satz: Welche Kosten- und Nutzenkomponenten werden gegenübergestellt? Welche Vanity-Metrik würden Sie aus dieser Liste \emph{nicht} als Hauptindikator wählen \textendash{} und warum?

\ytquery{Chatbot ROI KPI Bot Lösungsquote Eskalation Kosten pro Kontakt B2B}

\antwortlinien{12}

\clearpage

% =========================================================================
% L2
% =========================================================================
\section*{Niveau L2 \textendash{} Anwendung \& Transfer}
\addcontentsline{toc}{section}{L2 -- Anwendung}

% --- AUFGABE 4 -----------------------------------------------------------
\aufgabenkopf{4}{OfficeCloud Services \textendash{} Chatbot-Use-Cases}{\nivLii}{25\,min}

\begin{infobox}[Fallbeschreibung]
\textbf{Unternehmen:} ``OfficeCloud Services'' verkauft cloudbasierte Arbeitsplatzlösungen an kleine und mittlere Unternehmen. Viele Websitebesucher informieren sich über Preise, Funktionen, Datenschutz, Einrichtungsaufwand und Beratungsmöglichkeiten. Das Vertriebsteam erhält täglich viele ähnliche Fragen per E-Mail und Telefon.

\smallskip
\textbf{Gegeben:}
\begin{itemize}
  \item Websitebesucher stellen häufig Standardfragen.
  \item Viele Interessenten sind noch nicht kaufbereit.
  \item Vertriebsteam möchte mehr qualifizierte Beratungstermine.
  \item Kunden erwarten schnelle Antworten.
  \item Datenschutz und Datensicherheit sind wichtige Entscheidungskriterien.
  \item Budget für einen ersten Chatbot-Piloten ist begrenzt.
\end{itemize}
\end{infobox}

\noindent\textbf{Arbeitsauftrag:}

\begin{enumerate}
  \item Benennen Sie mindestens \textbf{fünf typische Kundenfragen}, die ein Chatbot beantworten könnte.
  \item Unterscheiden Sie diese in \emph{einfache Informationsfragen} und \emph{vertriebsrelevante Qualifizierungsfragen}.
  \item Beschreiben Sie, an welchen Punkten der Chatbot \textbf{an einen Menschen übergeben} sollte.
  \item Entwickeln Sie einen einfachen Dialogablauf in 6\,--\,8 Schritten: vom Websitebesuch bis zur Terminbuchung.
  \item Benennen Sie \textbf{drei Risiken}, die bei einem Chatbot-Einsatz entstehen könnten \textendash{} und einen Mitigation-Mechanismus je Risiko.
\end{enumerate}

\ytquery{B2B Chatbot Use Case Lead Qualifizierung Terminbuchung Dialogablauf}

\antwortlinien{22}

\clearpage

% --- AUFGABE 5 -----------------------------------------------------------
\aufgabenkopf{5}{OfficeCloud Services \textendash{} Tool-Bewertung und ROI}{\nivLii}{30\,min}

\begin{infobox}[Mögliche Tools]
\textbf{A:} Website-Chatbot für FAQ und Leadqualifizierung.\\
\textbf{B:} Terminbuchungstool mit Kalenderintegration.\\
\textbf{C:} Automatisierte E-Mail-Strecke nach Chatbot-Kontakt.\\
\textbf{D:} CRM-Integration zur Übergabe qualifizierter Leads.\\
\textbf{E:} KI-basierte Produktempfehlungen im Chat.\\
\textbf{F:} Dashboard zur Messung von Chatbot-Leads und Conversion.\\[4pt]
\textbf{Rahmen:} Budget 65.000\,\euro{} / 6 Monate \textperiodcentered{} Vertriebsteam stark ausgelastet \textperiodcentered{} IT-Ressourcen begrenzt \textperiodcentered{} Geschäftsführung erwartet messbare Ergebnisse nach 3 Monaten \textperiodcentered{} Kunden dürfen nicht frustriert werden.
\end{infobox}

\noindent\textbf{Arbeitsauftrag:}

\begin{enumerate}
  \item Bewerten Sie jedes Tool in fünf Dimensionen (Skala 1\,--\,5): \emph{Kundennutzen}, \emph{Aufwand}, \emph{Kosten}, \emph{Datenabhängigkeit}, \emph{Vertriebswirkung}.
  \item Priorisieren Sie die Tools in: \fachbegriff{sofort starten}, \fachbegriff{später testen}, \fachbegriff{zunächst nicht umsetzen}.
  \item Entscheiden Sie, welcher Use Case für den \emph{ersten} Pilot am sinnvollsten ist.
  \item Legen Sie \textbf{fünf KPIs} fest, an denen der Erfolg gemessen wird.
  \item Formulieren Sie eine einfache \textbf{ROI-Logik}: Welche Kosten entstehen, welcher Nutzen soll messbar werden, ab wann ist der Pilot wirtschaftlich?
\end{enumerate}

\medskip
\noindent\begin{tabularx}{\linewidth}{@{}l|c|c|c|c|c@{}}
\toprule
\textbf{Tool} & \textbf{Nutzen} & \textbf{Aufwand} & \textbf{Kosten} & \textbf{Daten} & \textbf{Vertrieb} \\
\midrule
A: FAQ-Bot          & \rule{1.2em}{0.4pt} & \rule{1.2em}{0.4pt} & \rule{1.2em}{0.4pt} & \rule{1.2em}{0.4pt} & \rule{1.2em}{0.4pt} \\
B: Termin-Buchung   & \rule{1.2em}{0.4pt} & \rule{1.2em}{0.4pt} & \rule{1.2em}{0.4pt} & \rule{1.2em}{0.4pt} & \rule{1.2em}{0.4pt} \\
C: E-Mail-Strecke   & \rule{1.2em}{0.4pt} & \rule{1.2em}{0.4pt} & \rule{1.2em}{0.4pt} & \rule{1.2em}{0.4pt} & \rule{1.2em}{0.4pt} \\
D: CRM-Integration  & \rule{1.2em}{0.4pt} & \rule{1.2em}{0.4pt} & \rule{1.2em}{0.4pt} & \rule{1.2em}{0.4pt} & \rule{1.2em}{0.4pt} \\
E: KI-Empfehlungen  & \rule{1.2em}{0.4pt} & \rule{1.2em}{0.4pt} & \rule{1.2em}{0.4pt} & \rule{1.2em}{0.4pt} & \rule{1.2em}{0.4pt} \\
F: Dashboard        & \rule{1.2em}{0.4pt} & \rule{1.2em}{0.4pt} & \rule{1.2em}{0.4pt} & \rule{1.2em}{0.4pt} & \rule{1.2em}{0.4pt} \\
\bottomrule
\end{tabularx}

\ytquery{Chatbot Implementierung Pilot ROI KPI Tool Bewertung B2B}

\antwortlinien{16}

\clearpage

% --- AUFGABE 6 -----------------------------------------------------------
\aufgabenkopf{6}{Dialog-Design \textendash{} Übergaberegeln Bot \textrightarrow{} Mensch}{\nivLii}{20\,min}

\noindent\textbf{Arbeitsauftrag:}

\noindent Entwerfen Sie eine \textbf{Eskalationslogik} für einen B2B-Chatbot. Beantworten Sie für die folgenden sechs Situationen:

\begin{enumerate}
  \item Kunde fragt dreimal hintereinander die gleiche Frage, weil die Antwort nicht passt.
  \item Kunde tippt: ``Ich möchte mit einem Menschen sprechen''.
  \item Kunde beschreibt einen Sonderfall, der nicht in den Standardpfaden vorgesehen ist.
  \item Kunde erwähnt rechtliche oder Compliance-Themen (Datenschutz, Vertragsausstieg).
  \item Kunde äußert Frust oder Beschwerde.
  \item Kunde fragt nach einem konkreten Angebot $>$\,10.000\,\euro{}.
\end{enumerate}

\noindent Für jede Situation:
\begin{itemize}
  \item Eskalationsregel (z.\,B.\ ``nach 3 wiederholten Fragen'', ``bei Schlüsselwörtern wie \texttt{Beschwerde, Vertrag, Anwalt, Reklamation}'').
  \item An wen wird übergeben (Innendienst, Vertrieb, Servicemanagement, Rechts-/Compliance)?
  \item Welche Kontextinformation muss mit übergeben werden, damit der Mensch ohne Doppelfragen weiterhelfen kann?
\end{itemize}

\noindent Schließen Sie mit \emph{einem} Satz: Welche Eskalation darf der Bot \emph{nicht} verzögern?

\ytquery{Chatbot Eskalation Bot zu Mensch Übergabe B2B Dialog Design Regeln}

\antwortlinien{20}

\clearpage

% --- AUFGABE 7 -----------------------------------------------------------
\aufgabenkopf{7}{B2B Supply Solutions \textendash{} Use-Case-Priorisierung}{\nivLii}{22\,min}

\begin{infobox}[Kontext]
``B2B Supply Solutions'' verkauft technische Verbrauchsmaterialien, Ersatzteile und Servicepakete an mittelständische Industriekunden. Über Website, Kundenportal und E-Mail kommen monatlich ca.\ \textbf{4.000 Anfragen}, davon ca.\ \textbf{45\,\%} reine Standardinformationen (Verfügbarkeit, Lieferzeit, Preise, Kompatibilität).

\smallskip
\textbf{Ziele 9 Monate:} 25\,\% weniger Standardanfragen im Innendienst \textperiodcentered{} 20\,\% schnellere Reaktionszeit \textperiodcentered{} 15\,\% mehr qualifizierte digitale Leads. Budget: 220.000\,\euro.
\end{infobox}

\noindent\textbf{Mögliche Use Cases:}

\begin{enumerate}
  \item \textbf{FAQ-Chatbot} für Standardfragen.
  \item \textbf{Produktverfügbarkeits- und Lieferzeit-Auskunft} (Integration ERP).
  \item \textbf{Kompatibilitäts-Assistent} (Eingabe Maschinen-/Geräte-Modell).
  \item \textbf{Nachbestellung über Chat} mit Login.
  \item \textbf{Leadqualifizierung} für Großanfragen.
  \item \textbf{Servicepaket-Empfehlung} bei wiederkehrenden Anfragen.
\end{enumerate}

\noindent\textbf{Arbeitsauftrag:}

\begin{enumerate}
  \item Bewerten Sie jeden Use Case nach: \emph{Volumen / Häufigkeit}, \emph{Automatisierbarkeit (klar regelbasiert vs.\ erklärungsbedürftig)}, \emph{Datenabhängigkeit (ERP / CRM / Login)}, \emph{Kundennutzen}.
  \item Wählen Sie \textbf{zwei Use Cases} für die ersten 3 Monate \textendash{} mit der besten Wirkung-zu-Aufwand-Bilanz.
  \item Wählen Sie \textbf{einen Use Case}, den Sie bewusst \emph{nicht} im ersten Schritt umsetzen \textendash{} obwohl er attraktiv erscheint.
  \item Begründen Sie das Schätzgerüst: Wenn 45\,\% der 4.000 Anfragen (\,= 1.800) automatisierbar sind und der Bot davon 60\,\% löst, entlastet das wie viele Innendienstkapazitäten pro Monat (Annahme: 12\,min pro Anfrage)?
\end{enumerate}

\ytquery{B2B Chatbot Use Case Priorisierung Industrie Ersatzteile Anfragen Innendienst}

\antwortlinien{20}

\clearpage

% --- AUFGABE 8 -----------------------------------------------------------
\aufgabenkopf{8}{Wenn der Bot der Marke schadet}{\nivLii}{20\,min}

\noindent\textbf{Arbeitsauftrag:}

\noindent Diskutieren Sie drei realistische Szenarien, in denen ein Chatbot \emph{Schaden} anrichtet. Beschreiben Sie für jedes Szenario:

\begin{enumerate}
  \item \textbf{Szenario A:} Der Bot beantwortet eine Datenschutzfrage falsch (Verarbeitung von Kundendaten, EU-Speicherung).
  \item \textbf{Szenario B:} Der Bot reagiert auf eine emotionale Reklamation mit einem höflichen Standardtext.
  \item \textbf{Szenario C:} Der Bot empfiehlt KI-basiert ein Produkt, das das Unternehmen \emph{nicht} mehr im Sortiment hat.
\end{enumerate}

\noindent Für jedes Szenario:

\begin{itemize}
  \item Welcher konkrete Schaden entsteht (rechtlich, Reputation, Vertrieb, Compliance)?
  \item Welche \emph{technische} Schutzmaßnahme verhindert das Szenario?
  \item Welche \emph{organisatorische} Maßnahme (Verantwortlichkeit, Review-Prozess) muss zusätzlich greifen?
\end{itemize}

\noindent Schließen Sie mit einer Empfehlung: Welche \emph{drei Themen} sollten in einem Chatbot \emph{nie} automatisiert beantwortet werden \textendash{} sondern immer mit klarer Eskalation an einen Menschen?

\ytquery{Chatbot Risiken Marke Datenschutz Reklamation Compliance B2B}

\antwortlinien{20}

\clearpage

% =========================================================================
% L3
% =========================================================================
\section*{Niveau L3 \textendash{} Management \& Entscheidungsarchitektur}
\addcontentsline{toc}{section}{L3 -- Management}

% --- AUFGABE 9 -----------------------------------------------------------
\aufgabenkopf{9}{Fallstudie B2B Supply Solutions \textendash{} Conversational-Commerce-Architektur}{\nivLiii}{45\,min}

\begin{infobox}[Fallstudie: B2B Supply Solutions]
``B2B Supply Solutions'' verkauft technische Verbrauchsmaterialien, Ersatzteile und Servicepakete. Anfragen kommen über Website, Kundenportal und E-Mail. Der Innendienst ist stark ausgelastet, viele Anfragen sind wiederkehrend.

\smallskip
\textbf{Rahmenbedingungen:}
\begin{itemize}
  \item Jahresumsatz: 42\,Mio.\,\euro
  \item Innendienst: 16 Personen \textperiodcentered{} Außendienst: 12 \textperiodcentered{} Service: variabel
  \item Websitebesuche pro Monat: 75.000 \textperiodcentered{} Portal-Nutzer: 8.500
  \item Anfragen pro Monat: 4.000 (45\,\% Standardinfo)
  \item Budget für 9 Monate: 220.000\,\euro
  \item Ziele: 25\,\% weniger Standardanfragen im Innendienst, 20\,\% schnellere Reaktion, 15\,\% mehr qualifizierte digitale Leads.
\end{itemize}

\smallskip
\textbf{Mögliche Use Cases:} FAQ-Chatbot \textperiodcentered{} Produktverfügbarkeit/Lieferzeit (ERP-integriert) \textperiodcentered{} Kompatibilitäts-Assistent \textperiodcentered{} Nachbestellung über Chat \textperiodcentered{} Leadqualifizierung Großanfragen \textperiodcentered{} Servicepaket-Empfehlung \textperiodcentered{} Reaktivierung Inaktive \textperiodcentered{} Statusabfragen.
\end{infobox}

\noindent\textbf{Arbeitsauftrag:}

\begin{enumerate}
  \item Analysieren Sie die Vertriebs- und Servicesituation in 3\,--\,5 Punkten.
  \item Identifizieren Sie die wichtigsten Engpässe und Hebel.
  \item Bewerten Sie alle Use Cases nach: Wirkung \textperiodcentered{} Aufwand \textperiodcentered{} Datenreife \textperiodcentered{} Akzeptanz Kunden \textperiodcentered{} Akzeptanz Vertrieb/Innendienst \textperiodcentered{} ROI-Potenzial.
  \item Entwickeln Sie eine \textbf{priorisierte 9-Monats-Roadmap} in drei Phasen (1\,--\,3, 4\,--\,6, 7\,--\,9).
  \item Definieren Sie die \textbf{Übergabe- und Eskalationslogik} (Bot \textrightarrow{} Innendienst, Bot \textrightarrow{} Service, Bot \textrightarrow{} Außendienst).
  \item Bauen Sie eine \textbf{KPI- und ROI-Architektur} auf (Nutzung, Lösungsquote, Eskalation, Conversion, Reaktionszeit, Kundenzufriedenheit, Kosten pro Kontakt, Umsatzbeitrag).
  \item Verteilen Sie das Budget (220.000\,\euro) grob auf Ihre Maßnahmen.
  \item Treffen Sie mindestens \emph{eine Entscheidung unter Unsicherheit} und begründen Sie diese aus Senior-Sicht.
\end{enumerate}

\ytquery{B2B Conversational Commerce Architektur Industriekunden Roadmap KPI ROI}

\antwortlinien{34}

\clearpage

% --- AUFGABE 10 ----------------------------------------------------------
\aufgabenkopf{10}{Grenzen der Automatisierung \textendash{} was bewusst menschlich bleibt}{\nivLiii}{30\,min}

\noindent\textbf{Diskussionsaufgabe.}

\noindent Die wirtschaftlich verlockendste Frage ist: ``Wie viel mehr können wir automatisieren?'' Die strategisch relevante Frage ist: ``Welche Kundeninteraktion \emph{darf} nicht automatisiert werden \textendash{} weil Vertrauen, Differenzierung oder Reputation sonst gefährdet sind?''

\medskip
\noindent\textbf{Arbeitsauftrag:}

\begin{enumerate}
  \item Beschreiben Sie in 3\,--\,4 Sätzen, warum \emph{Bot-Nutzungsrate} und \emph{Lösungsquote} allein keine ROI-Aussage liefern. Welche \emph{harten} und \emph{weichen} Kostenfaktoren gehören in eine seriöse ROI-Rechnung?
  \item Wählen Sie \textbf{drei Kundenkontakte}, die in einem mittelständischen B2B-Unternehmen \emph{bewusst menschlich} bleiben sollten \textendash{} und erklären Sie pro Kontakt:
  \begin{itemize}
    \item Warum dürfen Automatisierung und KI hier nicht eingesetzt werden?
    \item Welcher Markenschaden entsteht im \emph{worst case}, wenn doch automatisiert wird?
    \item Welche Governance-Regel verhindert, dass spätere Effizienz-Initiativen diese Grenze überschreiben?
  \end{itemize}
  \item Formulieren Sie eine \textbf{Faustregel} (1\,--\,2 Sätze): Ab wann ist eine zusätzliche Automatisierung wirtschaftlich \emph{und} verantwortbar \textendash{} und ab wann wird sie eine Markenfalle?
\end{enumerate}

\ytquery{Automatisierung Grenzen B2B Vertrauen Marke menschlich Senior Management}

\antwortlinien{20}

\clearpage

% --- AUFGABE 11 ----------------------------------------------------------
\aufgabenkopf{11}{Transferprojekt \textendash{} Conversational-Commerce-Architektur 12 Monate}{\nivLiii}{45\,min}

\begin{infobox}[Rolle und Ausgangslage]
\textbf{Ihre Rolle:} Chief Revenue Officer (oder Digital Sales Architect, Leiter:in Vertriebstransformation, Conversational-Commerce-Manager:in, CRM- und Automation-Leiter:in, externe:r Berater:in).

\smallskip
\textbf{Ausgangslage:} Ein mittelständisches Unternehmen möchte Chatbots und digitale Vertriebstools einführen, um Kundeninteraktion zu automatisieren, den Vertrieb zu entlasten und digitale Leads zu qualifizieren. Es gibt Website, CRM, Kundenportal, E-Mail-Marketing und erste digitale Kampagnen. Kundenprozesse sind uneinheitlich, Standardfragen belasten den Innendienst, Lead-Übergaben sind nicht sauber, Managemententscheidungen beruhen auf unvollständigen Daten.

\smallskip
\textbf{Rahmenbedingungen:}
\begin{itemize}
  \item Budget begrenzt \textperiodcentered{} Zeitdruck durch digitale Wettbewerber.
  \item Vertrieb und Service stark ausgelastet \textperiodcentered{} IT-Ressourcen begrenzt.
  \item Datenqualität uneinheitlich.
  \item Kunden erwarten schnelle Antworten, aber auch fachliche Sicherheit.
  \item Datenschutz und Einwilligungen müssen berücksichtigt werden.
  \item Geschäftsführung erwartet messbaren ROI \textendash{} bei unsicheren Annahmen.
\end{itemize}
\end{infobox}

\noindent\textbf{Arbeitsauftrag:} Entwickeln Sie eine strategische \emph{Entscheidungsarchitektur} \textendash{} keine Tool-Liste \textendash{} für Chatbots, Conversational Commerce und digitale Vertriebstools.

\begin{enumerate}
  \item \textbf{Zielbild} für automatisierte Kundeninteraktion im Vertrieb (3\,--\,5 Sätze).
  \item Customer-Journey-Analyse mit geeigneten Automatisierungspunkten und bewusst menschlichen Momenten.
  \item Priorisierung der Use Cases: FAQ, Leadqualifizierung, Terminbuchung, Produktberatung, Nachbestellung, Serviceanfragen, Reaktivierung.
  \item \textbf{Entscheidung}, welche Use Cases automatisiert, pilotiert, später geprüft oder bewusst nicht automatisiert werden.
  \item Prozessarchitektur mit \textbf{Übergaberegeln} an Vertrieb, Service, Innendienst, Management.
  \item Daten- und Systemarchitektur (CRM, Website, Portal, E-Mail, Reporting).
  \item \textbf{Governance-Struktur} für Bot-Inhalte, Datenschutz, Qualitätssicherung, Eskalation, Verantwortlichkeiten.
  \item \textbf{KPI- und ROI-Architektur}: Nutzung, Lösungsquote, Leadqualität, Conversion, Prozessentlastung, Zufriedenheit, Umsatzbeitrag.
  \item Roadmap für 12 Monate (Phasen).
  \item Drei \textbf{Managemententscheidungen unter Unsicherheit}.
\end{enumerate}

\noindent\textbf{Zusatzanforderung:} Treffen Sie mindestens eine bewusste \emph{Entscheidung gegen} eine attraktive Automatisierung \textendash{} weil Datenqualität, Kundenerlebnis, rechtliche Sicherheit oder Markenwirkung das nicht zulassen. Begründen Sie aus Senior-Level-Perspektive.

\ytquery{Chief Revenue Officer Conversational Commerce Chatbot Architektur 12 Monate}

\antwortlinien{38}

\clearpage

% --- Abschluss -----------------------------------------------------------
\section*{Abschluss}
\thispagestyle{plain}

\noindent Sie haben alle elf Aufgaben bearbeitet. Vergleichen Sie nun Ihre Antworten mit dem \emph{Lösungsheft Tag 14}.

\medskip
\begin{infobox}[Was Sie jetzt können sollten]
\begin{itemize}
  \item Chatbot-Typen und Conversational-Commerce-Einsatzfelder sauber abgrenzen.
  \item KPIs und eine ROI-Logik für digitale Vertriebstools entwerfen.
  \item Use Cases nach Wirkung und Datenreife priorisieren.
  \item Eskalations- und Übergaberegeln zwischen Bot und Mensch designen.
  \item Risiken für Marke und Compliance erkennen und mit Maßnahmen begegnen.
  \item Eine Roadmap mit Pilot, Skalierung und bewussten Grenzen entwickeln.
  \item Eine Entscheidungsarchitektur formulieren, die Effizienz, Kundennutzen, Datenschutz und Marke ausbalanciert.
\end{itemize}
\end{infobox}

\vspace{1cm}
\noindent{\color{lpAccent}\rule{\linewidth}{0.6pt}}\\[4pt]
{\footnotesize\sffamily\color{lpGray}Lernplatt \textperiodcentered{} by Can Siebert \textperiodcentered{} Kurs 22 (Bo 143) \textperiodcentered{} Tag 14 \textperiodcentered{} Aufgabenheft \textperiodcentered{} \textcopyright{} 2026}

\end{document}
